\documentclass[11pt,a4paper]{article}

% --- Packages ---
\usepackage[utf8]{inputenc}
\usepackage[T1]{fontenc}
\usepackage{geometry}
\geometry{margin=2.5cm}
\usepackage{times}
\usepackage{graphicx}
\usepackage{booktabs}
\usepackage{hyperref}
\usepackage{natbib}
\bibliographystyle{plainnat}

% --- Title ---
\title{LLM-Based Code Generation: A Mini Review of Three Landmark Approaches}
\author{Your Name\\
\small Your Affiliation\\
\small \texttt{your.email@example.com}}
\date{\today}

% =============================================================================
\begin{document}
\maketitle

% --- Abstract ---
\begin{abstract}
Large language models (LLMs) have transformed automatic code generation,
achieving performance that ranges from solving introductory programming problems
to competing in programming contests.
This mini review compares three landmark approaches---Codex, AlphaCode, and
Code Llama---examining their architectures, training strategies, and evaluation
results.
We highlight the distinct design philosophies behind each system and discuss
implications for future research.
\end{abstract}

% --- 1. Introduction ---
\section{Introduction}
\label{sec:intro}

% TODO: Write introduction based on outline Section 1.
% Key points:
% - Emergence of LLMs for code generation
% - From traditional program synthesis to neural approaches
% - Scope: three representative systems
% - Research question

% --- 2. Background ---
\section{Background}
\label{sec:background}

% TODO: Write background based on outline Section 2.
% Key points:
% - Task definition (input: NL description, output: code)
% - Evaluation: pass@k, HumanEval, MBPP, Codeforces
% - Brief pre-LLM context

% --- 3. Approaches ---
\section{Approaches}
\label{sec:approaches}

\subsection{Codex}
\label{sec:codex}

% TODO: Write based on outline Section 3.1 + paper analysis.

\subsection{AlphaCode}
\label{sec:alphacode}

% TODO: Write based on outline Section 3.2 + paper analysis.

\subsection{Code Llama}
\label{sec:codellama}

% TODO: Write based on outline Section 3.3 + paper analysis.

% --- 4. Comparative Analysis ---
\section{Comparative Analysis}
\label{sec:comparison}

% TODO: Insert figure from outline Figure plan.
% \begin{figure}[htbp]
% \centering
% \includegraphics[width=0.9\textwidth]{figures/fig1.pdf}
% \caption{TODO: Caption from outline.}
% \label{fig:comparison}
% \end{figure}

% TODO: Insert comparison table from secondary/paper_comparison.md
% TODO: Write discussion based on outline Section 4.

% Placeholder table structure:
% \begin{table}[htbp]
% \centering
% \caption{Comparison of three LLM-based code generation approaches.}
% \label{tab:comparison}
% \begin{tabular}{lccc}
% \toprule
% \textbf{Aspect} & \textbf{Codex} & \textbf{AlphaCode} & \textbf{Code Llama} \\
% \midrule
% ... & ... & ... & ... \\
% \bottomrule
% \end{tabular}
% \end{table}

% --- 5. Conclusion ---
\section{Conclusion}
\label{sec:conclusion}

% TODO: Write conclusion based on outline Section 5.

% --- References ---
\bibliography{references}

\end{document}
